\documentclass[twocolumn,aps,prfluids]{revtex4-2}
\usepackage{amsmath,amssymb,graphicx}
\usepackage{hyperref}
\newcommand{\phig}{\varphi}

\begin{document}

\title{Golden-ratio scale invariance in Fibonacci-layered vortex circulation}

\author{Ziyavutdinov Magomed Kamalovich}
\affiliation{Independent researcher}
\email{zimakam@gmail.com}
\affiliation{ORCID: 0009-0005-9212-9921}

\date{\today}

\begin{abstract}
We show that a vortex whose circulation is organised as a Fibonacci hierarchy of K concentric layers exhibits an emergent phi-scale invariance in its core. Defining the local homogeneity exponent alpha(r) = -d log Gamma / d log r, we derive analytically that at the core scale r = r_c it satisfies alpha(K) = (K-1) log(phi) / 2 + c_inf + O(1/K), with c_inf approx -0.0665. The condition for exact self-similarity, Gamma(phi r) / Gamma(r) = 1/phi, requires alpha = 1 and therefore selects K* = 5.43, whose unique integer realisation is K = 5. Numerical evaluation for K in [2, 23] confirms the analytic formula to better than 1 percent. We discuss the geometric origin of the selected hierarchy and its relation to the regularised VORTEX profile used in SPARC rotation-curve fitting.
\end{abstract}

\maketitle
\section{Introduction}
The golden ratio phi = (1+sqrt(5))/2 appears across scales in nature, from phyllotaxis to quantum spectra. In fluid dynamics, self-similar vortex structures are typically associated with power-law profiles. Here we show that a Fibonacci-weighted hierarchy of concentric vortex layers produces an emergent phi-self-similarity that uniquely selects K = 5 layers.

\section{Model: FibonacciVortex}
Let Gamma(r) denote the circulation of a vortex whose core is organised as K concentric layers weighted by the Fibonacci sequence F_k and modulated by phi:
\begin{equation}
\Gamma(r) = \Gamma_0 \sum_{k=0}^{K-1} F_k \, \phig^{-k r/r_c}.
\end{equation}
The generating function of the Fibonacci numbers has a pole at x = 1/phi, corresponding to r = r_c. This pole controls the behaviour of Gamma(r) in the core.

\section{Local homogeneity exponent}
Define alpha(r) = -d log Gamma / d log r. For the model above at r = r_c, a saddle-point evaluation using F_k approx phi^k / sqrt(5) gives
\begin{equation}
\alpha(K) = \frac{(K-1)\log\phig}{2} + c_\infty + O(1/K),
\end{equation}
with the universal constant c_inf approx -0.0665, independent of Gamma_0 and r_c.

\section{Selection of K = 5}
Exact phi-self-similarity requires Gamma(phi r)/Gamma(r) = 1/phi, equivalently alpha = 1. Solving for K:
\begin{equation}
K^* = 1 + \frac{2(1 - c_\infty)}{\log\phig} = 5.43.
\end{equation}
The unique integer realisation is K = 5.

\section{Numerical verification}
We computed alpha(r) numerically via central differences (h = 1e-5) for K in [2, 23] and r in [0.3, 3.0] r_c. Table values below.

\begin{table}[h]
\caption{Local homogeneity exponent alpha(r_c): numerical vs analytic.}
\begin{ruledtabular}
\begin{tabular}{rrrr}
K & alpha num & alpha ana & diff \\
\hline
 5 & +0.9056 & +0.8959 & +0.0097 \\
 8 & +1.6235 & +1.6177 & +0.0058 \\
12 & +2.5841 & +2.5802 & +0.0039 \\
15 & +3.3051 & +3.3020 & +0.0031 \\
20 & +4.5074 & +4.5050 & +0.0024 \\
23 & +5.2289 & +5.2268 & +0.0021 \\
\end{tabular}
\end{ruledtabular}
\end{table}

The profile alpha(r) for K = 5 peaks at alpha = 1.0005 near r = 1.5 r_c, confirming the core region of exact phi-self-similarity.

\section{Relation to SPARC VORTEX}
The regularised VORTEX profile used in SPARC rotation-curve fitting, v^2 = A^2 (x - arctan x) / x with x = r/r_c, behaves as v ~ 1/sqrt(r) at large r, in contrast to the 1/r behaviour here. Both share the property of a regularised core controlled by a single length scale r_c, suggesting a common geometric origin.

\section{Discussion}
The result K = 5 is robust: it follows from the condition alpha = 1 and does not require fine-tuning. Generalisations to other irrational bases lambda != phi yield different K*, opening a family of scale-invariant vortex hierarchies.

\section{Conclusion}
A Fibonacci-layered vortex with K = 5 layers realises exact phi-self-similarity in its core. The value K = 5 is not a free parameter but a consequence of the golden-ratio structure of the circulation.

\begin{acknowledgments}
The author thanks the SPARC collaboration for public rotation-curve data used in related work.
\end{acknowledgments}

\bibliography{bibliography}

\end{document}
